\documentclass[11pt,a4paper]{article}
\usepackage[utf8]{inputenc}
\usepackage[T1]{fontenc}
\usepackage{amsmath}
\usepackage{amssymb}
\usepackage{booktabs}
\usepackage{hyperref}
\usepackage{geometry}
\usepackage{listings}
\usepackage{xcolor}
\usepackage{textcomp} % for extra symbols

\geometry{margin=1in}

\definecolor{codegreen}{rgb}{0,0.6,0}
\definecolor{codegray}{rgb}{0.5,0.5,0.5}
\definecolor{codepurple}{rgb}{0.58,0,0.82}
\definecolor{backcolour}{rgb}{0.95,0.95,0.92}

\lstset{
    backgroundcolor=\color{backcolour},   
    commentstyle=\color{codegreen},
    keywordstyle=\color{magenta},
    numberstyle=\tiny\color{codegray},
    stringstyle=\color{codepurple},
    basicstyle=\ttfamily\small,
    breakatwhitespace=false,         
    breaklines=true,                 
    captionpos=b,                    
    keepspaces=true,                 
    numbers=left,                    
    numbersep=5pt,                  
    showspaces=false,                
    showstringspaces=false,
    showtabs=false,                  
    tabsize=2,
    language=Python,
    extendedchars=true,
    literate={ε}{{$\epsilon$}}1
}

\title{Vienna Routing: Cost and Heuristics Documentation}
\author{Vienna Traffic Router Team}
\date{\today}

\begin{document}

\maketitle

This document explains how distances are measured, how edge costs ($g(n)$) are calculated, and how heuristics ($h(n)$) are combined to guide search algorithms like A*.

\section{Distance Calculation}

All distance measurements in the Vienna project use the \textbf{Haversine Formula}. This calculates the great-circle distance between two points defined by their latitude and longitude.

\subsection{Formula}
The distance $d$ is calculated as follows:
\begin{align*}
a &= \sin^2\left(\frac{\Delta\phi}{2}\right) + \cos(\phi_1)\cos(\phi_2)\sin^2\left(\frac{\Delta\lambda}{2}\right) \\
c &= 2 \cdot \operatorname{atan2}\left(\sqrt{a}, \sqrt{1-a}\right) \\
d &= R \cdot c
\end{align*}
Where $R = 6,371,000$ meters.

\subsection{Implementation Source}
In \texttt{src/graph/builder.py}:
\begin{lstlisting}
def haversine(lat1: float, lon1: float, lat2: float, lon2: float) -> float:
    R = 6_371_000
    phi1, phi2 = math.radians(lat1), math.radians(lat2)
    dphi = math.radians(lat2 - lat1)
    dlam = math.radians(lon2 - lon1)
    a = math.sin(dphi / 2) ** 2 + math.cos(phi1) * math.cos(phi2) * math.sin(dlam / 2) ** 2
    return R * 2 * math.atan2(math.sqrt(a), math.sqrt(1 - a))
\end{lstlisting}

\section{Edge Cost ($g(n)$)}

The edge cost, or path cost from the start to the current node ($g(n)$), represents the actual weight assigned to an edge in the graph.

\subsection{Effective Cost Calculation}
\begin{enumerate}
    \item \textbf{Manual Traffic Overrides}:
    \begin{itemize}
        \item \textbf{Blocked}: If intensity $\ge 95$, the edge is considered impassable ($\infty$ cost).
        \item \textbf{Sub-blocked}: Multiplier $= 0.5 + (\text{intensity}/100.0) \cdot 2.5$.
    \end{itemize}
    \item \textbf{Road Quality Multiplier}:
    Adjusted based on road type and surface.
\end{enumerate}

\subsection{Implementation Source}
In \texttt{src/graph/loader.py} $\rightarrow$ \texttt{get\_effective\_edge\_cost()}:
\begin{lstlisting}
    if overrides_map:
        edge_key = f'{edge["from"]}_{edge["to"]}'
        if edge_key in overrides_map:
            intensity = overrides_map[edge_key]
            if intensity >= 95:
                return float("inf")
            cost *= 0.5 + (intensity / 100.0) * 2.5

    if params is not None and params.get("use_road_quality", True):
        from src.heuristics.road_quality import get_road_quality_multiplier
        cost *= get_road_quality_multiplier(edge, params.get("vehicle_type"))
\end{lstlisting}

\section{Heuristic Cost ($h(n)$)}

The heuristic cost ($h(n)$) is an estimate of the remaining cost from the current node to the goal.

\subsection{Composite Multiplier}
The final $h(n)$ is:
\[ h(n) = \text{Haversine}(n, \text{goal}) \cdot \text{Multiplier}_{\text{Total}} \]

\section{Admissibility and Capping}

\subsection{Admissibility}
A heuristic is \textbf{admissible} if it never overestimates the true cost. The \textbf{Haversine distance} is the base admissible heuristic. However, multiplying it by factors $> 1.0$ (like weather penalties) can make it \textbf{inadmissible}.

\subsection{The MAX\_HEURISTIC\_CAP}
To prevent the search from becoming purely greedy, the final heuristic is capped:
\[ h(n)_{\text{final}} = \min(h, h_{\text{base}} \cdot \text{MAX\_HEURISTIC\_CAP}) \]

The default cap is \textbf{50.0}. This ensures \textbf{bounded-suboptimality} ($\epsilon$-admissibility), meaning the resulting path is guaranteed to be within a known factor of the optimal path.

\subsection{How Capping Helps}
\begin{itemize}
    \item \textbf{Search Focus}: Prevents $h(n)$ from ballooning and dwarfing $g(n)$.
    \item \textbf{Greedy Prevention}: Ensures A* doesn't blindly follow low $h$ values.
    \item \textbf{Stability}: Maintains consistent performance when multiple multipliers are stacked.
\end{itemize}

\subsection{Implementation Source}
In \texttt{src/heuristics/combined.py}:
\begin{lstlisting}
MAX_HEURISTIC_CAP = float(os.getenv("MAX_HEURISTIC_CAP", "50.0"))
# ...
        h = h_base * mult
        # Cap total h against arbitrary blow-up, keep epsilon-admissibility
        return min(h, h_base * MAX_HEURISTIC_CAP)
\end{lstlisting}

\section{Distance and Time Results}

\subsection{Implementation Source}
In \texttt{src/algorithms/base.py}:
\begin{lstlisting}
def sum_path_distance(graph: dict, path: list[str]) -> float:
    # ... loops through path and sums graph["edges"][edge_idx]["distance_m"] ...

VEHICLE_SPEEDS = {
    "car":        8.33,   # ~30 km/h urban
    "motorcycle": 9.72,   # ~35 km/h
}
\end{lstlisting}

\end{document}
